\setcounter{page}{25}
\section{\S2. $K(\cat{A})$ is Triangulated}

Let \(\cat{A}\) be an abelian category. A \textit{complex} of objects of \(\cat{A}\) is a collection \(X^\bullet = (X^n)_{n\in\mathbb{Z}}\) of objects of \(\cat{A}\), together with morphisms \(d^n\colon X^n\to X^{n+1}\) such that \(d^{n+1}\circ d^n = 0\) for all \(n\in\mathbb{Z}\).

A \textit{morphism of complexes} \(f\colon X^\bullet\to Y^\bullet\) is a family of morphisms \(f^n\colon X^n\to Y^n\) commuting with differentials:
\[
f^{n+1}\circ d_X^n = d_Y^n\circ f^n
\]
for all \(n\).

Two morphisms \(f,g\colon X^\bullet\to Y^\bullet\) are \textit{homotopic} if there exists a family of morphisms \(k=(k^n)\), \(k^n\colon X^n\to X^{n-1}\), such that
\[
f^n - g^n = d_Y^{n-1}\circ k^n + k^{n+1}\circ d_X^n
\]
for all \(n\). Homotopy is an equivalence relation, and compositions of homotopic maps are homotopic.

We define \(K(\cat{A})\) to be the category whose objects are complexes over \(\cat{A}\), and whose morphisms are homotopy equivalence classes of complex morphisms.

A complex \(X^\bullet\) is \textit{bounded below} if \(X^n=0\) for all sufficiently small \(n\ll 0\). Denote by \(K^+(\cat{A})\) the full subcategory of \(K(\cat{A})\) consisting of bounded-below complexes. Similarly define \(K^-(\cat{A})\) (bounded above) and \(K^b(\cat{A})\) (bounded on both sides).

\setcounter{page}{26}
Now we equip \(K(\cat{A})\) (and likewise \(K^+(\cat{A}),K^-(\cat{A}),K^b(\cat{A})\)) with a triangulated structure.
The translation functor \(T\) shifts degrees left by one and flips the sign of the differential:
\[
T(X^\bullet)^p = X^{p+1},\qquad d_{T(X)} = -d_X.
\]
We also write \(X^\bullet[n] = T^n(X^\bullet)\).

Given any morphism \(u\colon X^\bullet\to Y^\bullet\) in \(K(\cat{A})\), form its \textit{mapping cone} \(Z^\bullet\), which is the complex \(T(X^\bullet)\oplus Y^\bullet\), with differential given by the matrix
\[
\begin{pmatrix}
T(d_X) & T(u)\\
0 & d_Y
\end{pmatrix}.
\]
There are natural morphisms \(v\colon Y^\bullet\to Z^\bullet\) and \(w\colon Z^\bullet\to T(X^\bullet)\).

A \textit{triangle} in \(K(\cat{A})\) is any sextuple isomorphic to one obtained from the above construction starting from a morphism \(u\colon X^\bullet\to Y^\bullet\).

All axioms \(\text{TR1}\)--\(\text{TR4}\) are readily verified, using the fact that the mapping cone of the identity morphism \(\id_X\colon X^\bullet\to X^\bullet\) is homotopic to zero.
This cone is \(T(X^\bullet)\oplus X^\bullet\), and the matrix

\setcounter{page}{27}
\[
k = \begin{pmatrix}
0 & 0\\
\id_X & 0
\end{pmatrix}
\]
defines a homotopy to the zero complex.

\medskip
\textbf{Definition.}
Define the cohomology functor \(H\colon K(\cat{A})\to\cat{A}\) sending a complex \(X^\bullet\) to its \(0\)-th cohomology object:
\[
H(X^\bullet) = \ker d^0 \big/ \operatorname{im} d^{-1}.
\]
This is indeed a functor, since homotopic complex morphisms induce identical maps on cohomology.
For any \(i\in\mathbb{Z}\), write
\[
H^i = H\circ T^i.
\]

The functor \(H\) is a cohomological functor from \(K(\cat{A})\) to the abelian category \(\cat{A}\). It suffices to verify the long exact sequence for triangles coming from mapping cones of morphisms \(u\colon X^\bullet\to Y^\bullet\), where exactness holds by direct inspection.